\subsection{Quasi-deterministic finite-state automata}

\begin{definition}
	A quasi-deterministic $k$-separated finite-state automaton $\mathcal{A}$
	is a tuple
	\begin{align*}
		\langle \Sigma, \vec{Q}, X, s, \vec{F}, \vec{\tau}, \vec{\rho},
		\vec{\eta}, \eta_x \rangle
	\end{align*}
	where $k \in \mathbb{N}$, $k > 1$, $\Sigma$ is an alphabet, $Q_i$ is a set
	of quasi-states for $i \in \interval{1}{k}$, $X$ is a non-empty separator set
	such that $\forall i \in \interval{1}{k} (X \subseteq Q_i)$ and $\forall i,
	j \in \interval{1}{k}(i \neq j \implies Q_i \cap Q_j = X)$, $s \in
	\prod_{i=1}^k{Q_i}$ is a start state, $F_i \subseteq Q_i$ is a set of final
	quasi-states for $i \in \interval{1}{k}$, $\tau_i : Q_i \times \Sigma
	\rightarrow Q_i$ is a component quasi-transition function for $i \in
	\interval{1}{k}$, $\rho_i : Q_i \rightarrow X$ is a component reduction
	function for $i \in \interval{1}{k}$, $\eta_i : Q_i \times X \rightarrow
	Q_i$ is a component extension function for $i \in \interval{1}{k}$ and
	$\eta_x : X^k \rightarrow X$ is a separator extension function.
\end{definition}

\begin{remark}
	Let $\mathcal{A} = \langle \Sigma, \vec{Q}, X, s, \vec{F}, \vec{\tau},
	\vec{\rho}, \vec{\eta}, \eta_x \rangle$ be a quasi-deterministic
	$k$-separated finite-state automaton. We use $q^j$ for $j \in
	\interval{1}{k}$ to denote the $j$-th element of $q = \langle q^1, q^2,
	\ldots, q^k \rangle \in \prod_{i=1}^k{Q_i}$.
\end{remark}

\begin{definition}
	Let $\mathcal{A} = \langle \Sigma, \vec{Q}, X, s, \vec{F}, \vec{\tau},
	\vec{\rho}, \vec{\eta}, \eta_x \rangle$ be a quasi-deterministic
	$k$-separated finite-state automaton. The separator tansition function
	$\delta_x : \prod_{i=1}^k{Q_i} \times \Sigma \rightarrow X$ is defined as
	\begin{align*}
		\delta_x(q, \sigma) = \eta_x(\rho_1(\tau_1(q^1, \sigma)),
		\rho_2(\tau_2(q^2, \sigma)), \ldots, \rho_k(\tau_k(q^k, \sigma))).
	\end{align*}
\end{definition}

\begin{definition}
	Let $\mathcal{A} = \langle \Sigma, \vec{Q}, X, s, \vec{F}, \vec{\tau},
	\vec{\rho}, \vec{\eta}, \eta_x \rangle$ be a quasi-deterministic
	$k$-separated finite-state automaton. The component tansition functions
	$\delta_i : \prod_{j=1}^k{Q_j} \times \Sigma \rightarrow Q_i$ are defined
	as
	\begin{align*}
		\delta_i(q, \sigma) = \eta_i(\tau_i(q^i, \sigma), \delta_x(q,
		\sigma))
	\end{align*}
	for $i \in \interval{1}{k}$.
\end{definition}

\begin{definition}
	Let $\mathcal{A} = \langle \Sigma, \vec{Q}, X, s, \vec{F}, \vec{\tau},
	\vec{\rho}, \vec{\eta}, \eta_x \rangle$ be a quasi-deterministic
	$k$-separated finite-state automaton. The tansition function $\delta :
	\prod_{i=1}^k{Q_i} \times \Sigma \rightarrow \prod_{i=1}^k{Q_i}$ is defined
	as
	\begin{align*}
		\delta(q, \sigma) = \langle \delta_1(q, \sigma), \delta_2(q,
		\sigma), \ldots, \delta_k(q, \sigma) \rangle.
	\end{align*}
\end{definition}

\begin{definition}
	A path $\pi$ in a quasi-deterministic $k$-separated finite-state automaton
	$\mathcal{A} = \langle \Sigma, \vec{Q}, X, s, \vec{F}, \vec{\tau},
	\vec{\rho}, \vec{\eta}, \eta_x \rangle$ is a non-empty sequence of
	transitions
	\begin{align*}
		\langle q_1, \sigma_1, q_2 \rangle, \langle q_2, \sigma_2, q_3
		\rangle, \dots, \langle q_{n-1}, \sigma_{n-1}, q_n \rangle
	\end{align*}
	where $\forall i \in \interval{1}{n}(q_i \in \prod_{j=1}^k{Q_j})$, $\forall
	i \in \interval{1}{n-1} (\sigma_i \in \Sigma)$ and $\forall i \in
	\interval{1}{n-1} (q_{i+1} = \delta(q_i, \sigma_i))$. The path $\pi$ is
	successful if and only if $q_1 = s$ and $\exists i \in \interval{1}{k}
	(q_n^i \in F^i)$. The label of the path $\pi$ is the word
	$\sigma_1\sigma_2\ldots\sigma_{n-1}$ of length $n-1$ over $\Sigma$.
\end{definition}

\begin{definition}
	Let $\mathcal{A} = \langle \Sigma, \vec{Q}, X, s, \vec{F}, \vec{\tau},
	\vec{\rho}, \vec{\eta}, \eta_x \rangle$ be a quasi-deterministic
	$k$-separated finite-state automaton. The generalized function of
	transitions $\delta^* : \prod_{i=1}^k{Q_i} \times \Sigma^* \rightarrow
	\prod_{i=1}^k{Q_i}$ is defined inductively as the least function such that
	$\delta^*(q, \epsilon) = q$ for $q \in \prod_{i=1}^k{Q_i}$ and $\delta^*(q,
	w\sigma) = \delta(\delta^*(q, w), \sigma)$ for $q \in \prod_{i=1}^k{Q_i}$,
	$w \in \Sigma^*$ and $\sigma \in \Sigma$.
\end{definition}

\begin{definition}
	The language recognized by a quasi-deterministic $k$-separa\-ted finite-state
	automaton $\mathcal{A} = \langle \Sigma, \vec{Q}, X, s, \vec{F}, \vec{\tau},
	\vec{\rho}, \vec{\eta}, \eta_x \rangle$ is the set of all labels of
	successful paths
	\begin{align*}
		\mathcal{L}(\mathcal{A}) = \{ w \in \Sigma^* \mid \exists i \in
		\interval{1}{k}(\delta^*(s, w)^i \in F^i \}.
	\end{align*}
\end{definition}

\begin{construction}\label{conctruction_1}
	Let $\mathcal{A} = \langle \Sigma, \vec{Q}, X, S, F, \Delta \rangle$ be a
	$k$-separated finite-state automaton and let $P_i = Q_i \cup X$ for $i \in
	\interval{1}{k}$. We construct the quasi-deterministic $k$-separated
	finite-state automaton $\mathcal{A}'$ and the set of states $R$ inductively.

	\bigbreak
	Let $R^{(0)} = \emptyset$. At the base of the construction is a
	single-state quasi-deterministic $k$-separated finite-state automaton
	\begin{align*}
		\mathcal{A}'^{(1)} = \langle \Sigma, \vec{Q'^{(1)}}, 2^X, s,
		\vec{\emptyset}, \vec{\emptyset}, \vec{\emptyset}, \vec{\emptyset},
		\emptyset \rangle
	\end{align*}
	where
	\begin{itemize}
		\item $s = \langle \hat{C_\epsilon}(S) \cap P_1,
			\hat{C_\epsilon}(S) \cap P_2,
			\ldots, \hat{C_\epsilon}(S) \cap P_k \rangle$;
		\item $Q_i'^{(1)} = \{ s^i \}$ for $i \in \interval{1}{k}$;
		\item $R^{(1)} = \{ s \}$.
	\end{itemize}

	Assuming $\mathcal{A}'^{(n)} = \langle \Sigma, \vec{Q'^{(n)}}, 2^X, s,
	\vec{F'^{(n)}}, \vec{\tau^{(n)}}, \vec{\rho^{(n)}}, \vec{\eta^{(n)}},
	\eta_x^{(n)} \rangle$ and $R^{(n)}$ are defined we construct a
	quasi-deterministic $k$-separated finite-state automaton
	\begin{align*}
		\mathcal{A}'^{(n+1)} = \langle \Sigma,
		\vec{Q'^{(n+1)}}, 2^X, s, \vec{F'^{(n+1)}}, \vec{\tau^{(n+1)}},
		\vec{\rho^{(n+1)}}, \vec{\eta^{(n+1)}}, \eta_x^{(n+1)} \rangle
	\end{align*}
	where

	\begin{itemize}
		\item $F_i'^{(n+1)} = F_i'^{(n)} \cup \{ q^i \mid \exists q \in
			R^{(n)} \setminus R^{(n-1)}(q^i \cap F \neq \emptyset) \}$
			for $i \in \interval{1}{k}$;
		\item $\tau_i^{(n+1)} = \tau_i^{(n)} \cup \{ \langle q^i, \sigma,
			\hat{C_\epsilon}(\hat{\Delta}(q^i, \{ \sigma \})) \cap P_i
			\rangle \mid q \in R^{(n)} \setminus R^{(n-1)} \land \sigma
			\in \Sigma\}$ for $i \in \interval{1}{k}$;
		\item $\rho_i^{(n+1)} = \rho_i^{(n)} \cup \{ \langle q,  q \cap X
			\rangle \mid q \in Rng(\tau_i^{(n+1)}) \setminus
			Rng(\tau_i^{(n)}) \}$ for $i \in \interval{1}{k}$;
		\item $\eta_x^{(n+1)} = \eta_x^{(n)} \cup \{ \langle q,
			\bigcup_{i=1}^k \rho_i^{(n+1)}(\tau_i^{(n+1)}(q^i, \sigma))
			\rangle \mid q \in R^{(n)} \setminus R^{(n-1)} \land \sigma
			\in \Sigma \}$;
		\item $\eta_i^{(n+1)} = \eta_i^{(n)} \cup \{ \langle \tau_i^{(n+1)}
			(q^i, \sigma), \eta_x^{(n+1)}(q), \tau_i^{(n+1)}
			(q^i, \sigma) \cup \eta_x^{(n+1)}(q) \rangle
			\mid q \in R^{(n)} \setminus R^{(n-1)} \land \sigma \in
			\Sigma \}$ for $i \in \interval{1}{k}$;
		\item $Q_i'^{(n+1)} = Q_i'^{(n)} \cup Rng(\tau_i^{(n+1)}) \cup
			Rng(\eta_i^{(n+1)})$ for $i \in \interval{1}{k}$;
		\item $R^{(n+1)} = R^{(n)} \cup \{ \delta^{(n+1)}(q, \sigma) \mid
			\exists q \in R^{(n)} \setminus R^{(n-1)} \exists \sigma \in
			\Sigma \}$.
	\end{itemize}
\end{construction}

\begin{lemma}
	Let $\mathcal{A} = \langle \Sigma, \vec{Q}, X, S, F, \Delta \rangle$ be a
	$k$-separated finite-state automaton and let $P_i = Q_i \cup X$ for $i \in
	\interval{1}{k}$. Let $\mathcal{A}'^{(n)}$ be the quasi-deterministic
	$k$-separated finite-state automaton and let $R^{(n)}$ be the set of states
	derived from $\mathcal{A}$ by applying Construction \ref{conctruction_1}.
	There exists $n \leq |\prod_{i=1}^k{2^{P_i}}|$ such that $R^{(i)} =
	R^{(i+1)}$ and $\mathcal{A}^{(i)} = \mathcal{A}^{(i+1)}$ for $i \geq n$.
\end{lemma}

\begin{proof}
	For $n \in \mathbb{N}$ we have $R^{(n)} \subseteq R^{(n+1)} \subseteq
	\prod_{i=1}^k{2^{P_i}}$. Therefore, there exists $n \leq
	|\prod_{i=1}^k{2^{P_i}}|$ such that $R^{(n)} = R^{(n+1)}$. In that case,
	$\mathcal{A}'^{(n)} = \mathcal{A}'^{(n+1)}$. For the smallest such $n$ we
	declare $R = R^{(n)}$ and $\mathcal{A}' = \mathcal{A}'^{(n)}$.
\end{proof}

\begin{theorem}
	Let $\mathcal{A} = \langle \Sigma, \vec{Q}, X, S, F, \Delta \rangle$ be a
	$k$-separated finite-state automaton and let $P_i = Q_i \cup X$ for $i \in
	\interval{1}{k}$. Let $\mathcal{A}' = \langle \Sigma, \vec{Q'}, 2^X, s,
	\vec{F'}, \vec{\tau}, \vec{\rho}, \vec{\eta}, \eta_x \rangle$ be the
	quasi-deterministic $k$-separated finite-state automaton derived from
	$\mathcal{A}$ by applying Construction \ref{conctruction_1}. Then
	$\mathcal{L}(\mathcal{A'}) = \mathcal{L}(\mathcal{A})$.
\end{theorem}

\begin{proof}
	\bigbreak
	We show that every successful path in $\mathcal{A}$ is also a successful
	path in $\mathcal{A'}$ by induction on the length of the path.

	Let $\langle q_1, \sigma_1, q_2 \rangle, \langle q_2, \sigma_2, q_3 \rangle,
	\ldots, \langle q_{n-1}, \sigma_{n-1}, q_n \rangle$ be a successful path in
	$\mathcal{A}$. Let $I = \{ i_1, i_2, \ldots, i_l \} \subseteq \{ 1, 2,
	\ldots, n-1 \}$ be a set of indexes for $l \in \interval{1}{n-1}$ such that
	$\sigma_{i_j} \in \Sigma$ for $j \in \interval{1}{l}$ and $h < j$ if $i_h <
	i_j$ for $h, j \in \interval{1}{l}$.

	Let $p_1 = s$. We have $\sigma_{i_1} \in \Sigma$ and
	if $i_1 > 1$ then $\sigma_j = \epsilon$ for $j \in \interval{1}{i_1-1}$.
	Therefore, $q_{i_1} \in C_\epsilon(q_1) \subseteq \hat{C_\epsilon}(S) =
	\bigcup_{r=1}^k{\hat{C_\epsilon}(S) \cap P_r} = \bigcup_{r=1}^k{p_1^r}$.

	Let $p_h = \delta(p_1, \sigma_{i_1} \sigma_{i_2} \ldots
	\sigma_{i_h-1})$ and $q_{i_h} \in \bigcup_{r=1}^k{p_h^r}$ for
	$h \in \interval{2}{l}$.
	We define $p_{h+1} = \delta(p_h, \sigma_{i_h}) =
	\langle p_{h+1}^1, p_{h+1}^2, \ldots, p_{h+1}^k \rangle$ where
	\begin{align*}
		p_{h+1}^j &= \delta_j(p_h, \sigma_{i_h})\\
		&= \big[ \hat{C_\epsilon}(\hat{\Delta}(p_h^j, \{
		\sigma_{i_h} \})) \cap P_j \big] \cup
		\big[ \bigcup_{r = 1}^k{\hat{C_\epsilon}(\hat{\Delta}
		(p_h^r, \{ \sigma_{i_h} \})) \cap X} \big]\\
		&= \bigcup_{r = 1}^k{\hat{C_\epsilon}(\hat{\Delta}
		(p_h^r, \{ \sigma_{i_h} \})) \cap P_j}
	\end{align*}
	for $j \in \interval{1}{k}$.
	If $h < l$ then $\sigma_{i_{h+1}} \in \Sigma$ and if
	$i_{h+1} - i_h > 1$ then $\sigma_j = \epsilon$ for $j \in
	\interval{i_h+1}{i_{h+1}-1}$. Therefore,
	\begin{align*}
		q_{i_{h+1}} \in \hat{C_\epsilon}(\Delta(q_{i_h}, \sigma_{i_h}))
		&\subseteq \bigcup_{r=1}^k{\hat{C_\epsilon}(
		\hat{\Delta}(p_h^r, \{ \sigma_{i_h} \}))}\\
		&= \bigcup_{j=1}^k{\bigcup_{r=1}^k{
		\hat{C_\epsilon}(\hat{\Delta}(p_h^r, \{
		\sigma_{i_h} \})) \cap P_j}}
		= \bigcup_{j=1}^k{p_{h+1}^j}.
	\end{align*}
	Otherwise, $h = l$ and if $n - i_l > 2$ then $\sigma_j = \epsilon$ for $j
	\in \interval{i_l+1}{n-1}$. Therefore,
	\begin{align*}
		q_n \in \hat{C_\epsilon}(\Delta(q_{i_l}, \sigma_{i_l}))
		&\subseteq \bigcup_{r=1}^k{\hat{C_\epsilon}(
		\hat{\Delta}(p_l^r, \{ \sigma_{i_l} \}))}\\
		&= \bigcup_{j=1}^k{\bigcup_{r=1}^k{
		\hat{C_\epsilon}(\hat{\Delta}(p_l^r, \{
		\sigma_{i_l} \})) \cap P_j}}
		= \bigcup_{j=1}^k{p_{l+1}^j}.
	\end{align*}

	We have $q_n \in \bigcup_{j=1}^k{p_{l+1}^j}$ and $q_n \in F$. Therefore,
	$p_{l+1}^j \cap F \neq \emptyset$ and thus $p_{l+1}^j \in F_j'$ for some $j
	\in \interval{1}{k}$. Therefore, $\langle p_1, \sigma_{i_1}, p_2 \rangle$,
	$\langle p_2, \sigma_{i_2}, p_3 \rangle$, $\ldots$, $\langle p_l,
	\sigma_{i_l}, p_{l+1} \rangle$ is a successful path in $\mathcal{A'}$.

	Similarly, we show that every successful path in $\mathcal{A'}$ is also a
	successful path in $\mathcal{A}$.

	Let $\langle p_1, \sigma_1, p_2 \rangle, \langle p_2, \sigma_2, p_3 \rangle,
	\ldots, \langle p_{n-1}, \sigma_{n-1}, p_n \rangle$ be a successful path in
	$\mathcal{A'}$.
	There exists $j \in \interval{1}{k}$ such that $p_n^j \in F_j'$ and thus
	$p_n^j \cap F \neq \emptyset$. Therefore, there exists $q_n \in p_n^j$ such
	that $q_n \in F$.

	Let there be $j \in \interval{1}{k}$ and $q_i \in Q$ such that $q_i \in
	p_i^j$ for $i \in \interval{2}{n}$. We have
	\begin{align*}
		p_i^j &= \delta_j(p_{i-1}, \sigma_{i-1})\\
		&= \big[ \hat{C_\epsilon}(\hat{\Delta}(p_{i-1}^j, \{ \sigma_{i-1}
		\})) \cap P_j \big] \cup \big[ \bigcup_{r=1}^k{\hat{C_\epsilon}(
		\hat{\Delta} (p_{i-1}^r, \{ \sigma_{i-1} \})) \cap X} \big]\\
		&= \bigcup_{r=1}^k{\hat{C_\epsilon}(\hat{\Delta}(p_{i-1}^r, \{
		\sigma_{i-1} \})) \cap P_j}.
	\end{align*}
	Also, $q_i \in p_i^j \implies |p_i^j| > 0$. Then there exists $r \in
	\interval{1}{k}$ such that $q_i \in \hat{C_\epsilon}(\hat{\Delta}
	(p_{i-1}^r, \{ \sigma_{i-1} \})) \cap P_j$ and thus $|p_{i-1}^r| > 0$.
	Therefore, there exists $q_{i-1} \in Q$ such that $q_{i-1} \in p_{i-1}^r$.
	Moreover, since $q_i \in \hat{C_\epsilon}(\hat{\Delta}
	(p_{i-1}^r, \{ \sigma_{i-1} \}))$ there exists $l_{i-1} > 0$ and
	$q_{i-1, h} \in Q$ for $h \in \interval{1}{l_{i-1}}$ such that
	$\langle q_{i-1}, \sigma_{i-1}, q_{i-1, 1} \rangle \in \Delta$,
	$\langle q_{i-1, h}, \epsilon, q_{i-1, h+1} \rangle \in \Delta$ for $h \in
	\interval{1}{l_{i-1}-1}$ and $q_{i-1, l_{i-1}} = q_i$. Specifically, in
	case $l_{i-1} = 1$ we have $\langle q_{i-1}, \sigma_{i-1}, q_i \rangle \in
	\Delta$.

	There exists $j \in \interval{1}{k}$ and $q_1 \in Q$ such that $q_1 \in
	p_1^j$ and thus $q_1 \in S$. Therefore, $\langle q_1, \sigma_1, q_{1,1}
	\rangle$, $\langle q_{1,1}, \epsilon, q_{1,2} \rangle$, $\ldots$,
	$\langle q_{1,l_1-1}, \epsilon, q_2 \rangle$,
	$\langle q_2, \sigma_2, q_{2,1} \rangle$,
	$\langle q_{2,1}, \epsilon, q_{2,2} \rangle$,
	$\ldots$,
	$\langle q_{2,l_2-1}, \epsilon, q_3 \rangle$,
	$\ldots$,
	$\langle q_{n-1}, \sigma_{n-1}, q_{n-1,1} \rangle$,
	$\langle q_{n-1,1}, \epsilon, q_{n-1,2} \rangle$,
	$\ldots$,
	$\langle q_{n-1,l_{n-1}-1}, \epsilon, q_n \rangle$ is a
	successful path in $\mathcal{A}$.
\end{proof}
